%% ============================================================
%%  SOFTWAREX TEMPLATE - CLEAN DRAFT VERSION FOR WRITING
%%  Recommended for authoring/review. Switch to final layout later.
%% ============================================================

\documentclass[preprint,12pt]{elsarticle}

%% --- Packages ------------------------------------------------
\usepackage[utf8]{inputenc}
\usepackage[T1]{fontenc}
\usepackage{amsmath,amssymb}
\usepackage{graphicx}
\usepackage{booktabs}
\usepackage{array}
\usepackage{xcolor}
\usepackage{url}
\usepackage{hyperref}
\usepackage{listings}
\usepackage{microtype}
\usepackage[margin=2.5cm]{geometry}

%% --- Hyperref -----------------------------------------------
\hypersetup{
  colorlinks=true,
  linkcolor=blue,
  citecolor=blue,
  urlcolor=blue,
  pdfauthor={TODO: Authors},
  pdftitle={TODO: Software title / article title}
}

%% --- Listings -----------------------------------------------
\lstset{
  basicstyle=\ttfamily\small,
  breaklines=true,
  breakatwhitespace=false,
  columns=fullflexible,
  frame=single,
  inputencoding=utf8,
  extendedchars=true,
  showstringspaces=false,
  keepspaces=true,
  literate={á}{{\'a}}1 {é}{{\'e}}1 {í}{{\'i}}1 {ó}{{\'o}}1 {ú}{{\'u}}1
           {Á}{{\'A}}1 {É}{{\'E}}1 {Í}{{\'I}}1 {Ó}{{\'O}}1 {Ú}{{\'U}}1
           {ñ}{{\~n}}1 {Ñ}{{\~N}}1 {ü}{{\"u}}1 {Ü}{{\"U}}1
}

%% --- Bibliography style -------------------------------------
\bibliographystyle{elsarticle-num}

\begin{document}

\begin{frontmatter}

\title{TODO: Software title / article title}

\author[aff1]{TODO: Author One\corref{cor1}}
\ead{TODO: author.one@institution.edu}

\author[aff2]{TODO: Author Two}
\ead{TODO: author.two@institution.edu}

\cortext[cor1]{Corresponding author.}

\affiliation[aff1]{
  organization={TODO: Department, University/Institution},
  addressline={TODO: Address},
  city={TODO: City},
  postcode={TODO: Postal code},
  country={TODO: Country}
}

\affiliation[aff2]{
  organization={TODO: Department, University/Institution},
  addressline={TODO: Address},
  city={TODO: City},
  postcode={TODO: Postal code},
  country={TODO: Country}
}

\begin{abstract}
TODO: Write the abstract here. Maximum 250 words. Briefly describe the software,
its purpose, the problem it solves, its main functionalities, key results, and
the main contribution to the research community.
\end{abstract}

\begin{keyword}
TODO keyword 1 \sep
TODO keyword 2 \sep
TODO keyword 3 \sep
TODO keyword 4
\end{keyword}

\end{frontmatter}

\section*{Code metadata}
\label{sec:metadata}

\begin{table}[ht]
\centering
\small
\caption{Code metadata required by SoftwareX.}
\label{tab:metadata}
\begin{tabular}{p{0.08\linewidth}p{0.44\linewidth}p{0.40\linewidth}}
\toprule
\textbf{Nr.} & \textbf{Code metadata description} & \textbf{Metadata} \\
\midrule
C1 & Current code version & TODO: v1.0.0 \\
C2 & Permanent link to code/repository used for this code version & TODO: \url{https://github.com/user/repository} \\
C3 & Permanent link to reproducible capsule & TODO: \url{https://} or N/A \\
C4 & Legal Code License & TODO: MIT / Apache-2.0 / GPL-3.0 / other \\
C5 & Code versioning system used & TODO: git \\
C6 & Software code languages, tools, and services used & TODO: Python, R, C++, etc. \\
C7 & Compilation requirements, operating environments and dependencies & TODO: Python $\geq$ 3.9; dependencies: numpy, scipy, etc. \\
C8 & Link to developer documentation/manual & TODO: \url{https://} or N/A \\
C9 & Support email for questions & TODO: author.one@institution.edu \\
\bottomrule
\end{tabular}
\end{table}

\section{Motivation and significance}
\label{sec:motivation}

TODO: Explain why the software was developed. Describe the scientific or
computational problem it solves, why existing tools are insufficient, and how
the software fits into the research area.

\section{Software description}
\label{sec:description}

TODO: Provide a general description of the software, its scope, intended users,
main inputs, outputs, and expected usage context.

\subsection{Software architecture}
\label{sec:architecture}

TODO: Describe the software architecture. You may include a component diagram
as shown in Fig.~\ref{fig:architecture}.

\begin{figure}[ht]
  \centering
  %% Replace with your real figure:
  %% \includegraphics[width=0.82\linewidth]{figures/Figure_1.pdf}
  \fbox{\rule{0pt}{5cm}\rule{0.82\linewidth}{0pt}}
  \caption{TODO: Brief title of the figure. Provide a detailed explanation of
the architecture diagram, including all symbols and abbreviations used.}
  \label{fig:architecture}
\end{figure}

\subsection{Software functionalities}
\label{sec:functionalities}

TODO: Describe the main functionalities of the software.

\subsubsection{Functionality 1}
\label{sec:functionality1}

TODO: Describe functionality 1.

\subsubsection{Functionality 2}
\label{sec:functionality2}

TODO: Describe functionality 2.

\section{Illustrative examples}
\label{sec:examples}

TODO: Present one or more concrete use cases demonstrating the software. Include
code snippets, results, figures, or tables as appropriate.

\begin{table}[ht]
  \caption{TODO: Table title.}
  \label{tab:results}
  \centering
  \begin{tabular}{lcc}
    \toprule
    \textbf{Parameter} & \textbf{Value A} & \textbf{Value B} \\
    \midrule
    TODO row 1 & 0.00 & 0.00 \\
    TODO row 2 & 0.00 & 0.00 \\
    TODO row 3 & 0.00 & 0.00 \\
    \bottomrule
  \end{tabular}
\end{table}

\begin{equation}
  \label{eq:example}
  f(x) = \int_{-\infty}^{+\infty} \hat{f}(\xi)\, \exp(2\pi i \xi x)\, d\xi
\end{equation}

Equation~(\ref{eq:example}) shows TODO: explain what the equation represents.

\begin{lstlisting}[language=Python,
caption={TODO: Description of the code snippet.},
label={lst:example}]
# TODO: replace with your example code
import mymodule

result = mymodule.function(param1=value1,
                           param2=value2)
print(result)
\end{lstlisting}

\section{Impact}
\label{sec:impact}

TODO: Describe the expected or observed impact of the software on the scientific
community or practical applications. Explain who benefits from it and what new
research or applications it enables.

\section{Conclusions}
\label{sec:conclusions}

TODO: Summarize the main characteristics of the software and its contributions.
You may also mention limitations and future work.

\section*{CRediT authorship contribution statement}

\textbf{TODO: Author One:} Conceptualization, Software, Writing -- original draft.\\
\textbf{TODO: Author Two:} Methodology, Validation, Writing -- review \& editing.

\section*{Declaration of competing interest}

The authors declare that they have no known competing financial interests or
personal relationships that could have appeared to influence the work reported
in this paper.

\section*{Data availability}

The data that support the findings of this study are openly available in
TODO: repository name at \url{TODO: https://doi.org/xxxxx}.

\section*{Acknowledgements}

TODO: Acknowledge people who helped with language editing, writing, review, or
other support. If applicable, mention funding sources.

\bibliography{referencias}

\end{document}
